\section{\sphenix Detector}
\label{sec:detector}
\sphenix is designed to be a general-purpose detector, aiming to measure jet and heavy-flavor probes of the QGP created in \auau collisions at RHIC. A four-component precision tracking system enables measurements of heavy flavor and jet substructure observables, while the electromagnetic and hadronic calorimeter system is crucial for measuring the energy of jets and identifying direct photons and electrons. 
% Going outwards starting from the beam line, \sphenix comprises the following subsystems~\cite{Sphenix:TDR}: the MAPS-based Vertex Detector (\mvtx); the INTermediate Tracker (\intt)~\cite{Akiba:2025atb}; the Time Projection Chamber (\tpc)~\cite{Klest:2020sdb}; the Time Projection Chamber Outer Tracker (\tpot)~\cite{Aune:2024ysr}; the Electromagnetic Calorimeter (\emcal)~\cite{Aidala:2020toz,sPHENIX:2017lqb}; the Inner Hadronic Calorimeter (\ihcal)~\cite{sPHENIX:2017lqb};  the 1.4\,T superconducting solenoid magnet~\cite{OConnor:1998llb} and the Outer Hadronic Calorimeter (\ohcal)~\cite{sPHENIX:2017lqb}. 
% Except for \tpot, all detectors have full azimuthal coverage and span $\left|\eta\right| < 1.1$ in pseudorapidity. \sphenix also includes a number of forward detectors, namely the Minimum Bias Detectors (\mbd), the sPHENIX Event Plane Detectors (\sepd), and the Zero Degree Calorimeters (\zdc) which include the Shower Maximum Detector (\smd). 
%A 1.4~Tesla solenoid magnet separates the hadronic calorimeter into two segments.  
% \sphenix began its commissioning process in RHIC Run-2023 with Au+Au collisions. During RHIC Run-2024, \sphenix collected a large sample of $p$+$p$ physics data alongside a smaller sample of Au+Au data to complete its commissioning phase in that collision system. 
%Minimum bias events are triggered with the Minimum Bias Detector (MBD) which are located at $3.51 < |\eta | < 4.61$.
This analysis is based on data collected using the \intt, the Minimum Bias Detectors (\mbd), and the Zero Degree Calorimeters (\zdc), which are described in detail.

The \intt is a two-layer barrel strip tracker with a clamshell geometry~\cite{INTT_ladder_NIM_cite}. Each layer comprises two sensor rings, offset in $\phi$ to provide full azimuthal coverage, and is situated between the MAPS-based Vertex Detector (MVTX) and the Time Projection Chamber in \sphenix.
It is designed to provide the \sphenix tracking system with the capability to associate reconstructed tracks to individual RHIC bunch crossings, enabling effective out-of-time pileup discrimination and suppression. 
The detector consists of 56 silicon ladders, 24 in the inner and 32 in the outer barrel, evenly spaced in a cylindrical configuration around the beam pipe at radial positions of approximately 7.2, 7.8, 9.7, and 10.3\cm from the beam axis. The active area of each ladder measures 45.6\cm in length and 2\cm in width and is composed of four silicon sensors. This design ensures hermetic coverage over $2\pi$ in azimuth and $|\eta|<1.1$ in pseudorapidity for collision vertices within $\pm10\cm$ of the nominal interaction point along the beam axis. The silicon sensors of each ladder are segmented into 52 silicon blocks, with each block containing 128 strips. The strips have a pitch of 78\mum, and extend 16 or 20\mm along the beam axis, resulting in a total of 6,656 readout channels per ladder. To minimize material interference and preserve track reconstruction accuracy, each silicon ladder is designed with a radiation length \ensuremath{X/X_0} of 1.14\%. Before the \intt, toward the collision point, the material budget is approximately 3.8\% \ensuremath{X/X_0} due to contributions from the beam pipe (1.8\% \ensuremath{X/X_0}) and the MVTX detector (2.0\% \ensuremath{X/X_0}). The single-hit detection efficiency of the \intt ladder is measured to be better than 99\% in test beam. Table~\ref{tab:INTT_spec} summarizes key specifications of the \intt.

%In the data used in this analysis, the \intt was operated in a trigger-plus-extended readout mode, taking collision data from a 56$\times$56 bunch RHIC store, which is half of the nominal capacity of RHIC. The trigger-plus-extended readout mode extends the hit acceptance window to 100 RHIC beam clock oscillation (BCO) after a triggered collision associated with a Global-Level 1 (GL1) trigger signal, effectively capturing collision hits from subsequent bunch crossings, as shown in Figure~\ref{fig:INTT_timing}. The pattern of evenly-spaced spikes observed correspond to the 56 RHIC bunch crossings, demonstrating the \intt's single-bunch-crossing timing resolution. 

% In the data used for this analysis, the \intt was operated in a trigger-plus-extended readout mode, recording collision data from a 56$\times$56 bunch RHIC store, where the bunches were evenly spaced within the beam pipes. Each RHIC ring can accommodate up to 111 beam bunches, resulting in empty crossings between filled bunch crossings in this RHIC store. Figure~\ref{fig:INTT_timing} shows the INTT hit counts as a function of beam crossing. The trigger-plus-extended readout mode extends the hit-acceptance window following a triggered collision associated with a Global-Level 1 (GL1) trigger signal (red histogram), allowing data acquisition over 100 bunch crossings (BCO) in the nominal RHIC configuration. The observed pattern of evenly spaced spikes (blue histogram) corresponds to hits collected from the rest of 56 RHIC bunch crossings, demonstrating the \intt's single-bunch-crossing timing resolution.

\renewcommand{\arraystretch}{1.3}
\begin{table}[t!]
    \centering
    \caption{Specification of the \sphenix Intermediate Silicon Tracker (\intt).}
    \begin{tabular}{l c}
        \toprule[1pt]
        \textbf{Element} & \textbf{Value} \\
        \midrule[1pt]
        Number of ladders & 56 \\
        \multirow{2}{*}{Barrel radial distance to the beam line}  & 7.2 and 7.8\cm (Inner layer) \\
         & 9.7 and 10.3\cm (Outer layer)\\
        Radiation length (per ladder) & 1.14\% [\ensuremath{X/X_0}] \\
        Active area (per ladder) & $45.6 \times 2 \cma$  \\
        Number of channels (per ladder) & 6,656 \\
        Channel strip pitch & 78 \mum \\
        \multirow{2}{*}{Channel strip length} & 16\mm (strips within $\pm$13.0\cm along the beam axis) \\
         & 20\mm (otherwise) \\
        \bottomrule[1pt]
    \end{tabular}
    \label{tab:INTT_spec}
\end{table}
\renewcommand{\arraystretch}{1}

The \mbd is located on both sides of the interaction point at $\pm250\cm$ along the beam axis, covering the pseudorapidity range $3.51 < |\eta| < 4.61$~\cite{PHENIX:2003tlh}. Each side consists of 64 photomultiplier tubes (PMTs). The PMTs are arranged in three concentric rings around the beam pipe, ensuring full $2\pi$ azimuthal coverage. The \mbd serves as the primary detector for triggering minimum-bias (MB) events in heavy-ion collisions and provides essential global event information, including the event centrality and the measurement of the collision z-vertex position.

The \zdc of sPHENIX is located at a distance of approximately 18\m from both sides of the interaction point~\cite{Adler:2000bd}. It is a longitudinally segmented sampling hadronic calorimeter constructed from a tungsten alloy and optical fibers. These fibers transmit Cherenkov light produced by secondary charged particles within hadronic showers to PMTs. The \zdc is calibrated so that the single-neutron peak is positioned at the nominal value of 100 GeV. The ZDC is used in the determination of the MB criteria to distinguish \mbd-triggered events from beam-induced background.